\documentclass[11pt,letterpaper,oneside]{report}
\usepackage[T1]{fontenc}
\usepackage[utf8]{inputenc}
\usepackage[english]{babel}
\usepackage{lmodern}
\usepackage{microtype}
\usepackage[margin=0.86in,headheight=16pt]{geometry}
\usepackage[table,dvipsnames]{xcolor}
\usepackage{graphicx}
\usepackage{booktabs,tabularx,longtable,array}
\usepackage{enumitem}
\usepackage{titlesec}
\usepackage{fancyhdr}
\usepackage{hyperref}
\usepackage{xurl}
\Urlmuskip=0mu plus 1mu\relax
\usepackage{tcolorbox}
\usepackage{listings}
\usepackage{tikz}
\usetikzlibrary{arrows.meta,positioning,shapes.geometric,fit,calc}
\usepackage{needspace}
\usepackage{amssymb}
\definecolor{ink}{HTML}{14253A}
\definecolor{navy}{HTML}{17385D}
\definecolor{blue}{HTML}{20739A}
\definecolor{pale}{HTML}{ECF4F8}
\definecolor{green}{HTML}{1D6C5C}
\definecolor{softgreen}{HTML}{EEF8F4}
\definecolor{border}{HTML}{D4DFE7}
\definecolor{grey}{HTML}{506174}
\hypersetup{colorlinks=true,linkcolor=navy,urlcolor=blue,pdfauthor={Micro-Frontend Engineering Documentation Suite}}
\titleformat{\chapter}[display]{\normalfont\sffamily\bfseries\color{ink}}{\Large\color{blue}\MakeUppercase{\chaptername} \thechapter}{8pt}{\Huge}
\titlespacing*{\chapter}{0pt}{-10pt}{24pt}
\titleformat{\section}{\Large\sffamily\bfseries\color{navy}}{\thesection}{0.7em}{}
\titleformat{\subsection}{\large\sffamily\bfseries\color{navy}}{\thesubsection}{0.6em}{}
\setlist[itemize]{itemsep=3pt,topsep=5pt,leftmargin=1.6em}
\setlist[enumerate]{itemsep=4pt,topsep=5pt,leftmargin=1.8em}
\pagestyle{fancy}\fancyhf{}
\fancyhead[L]{\footnotesize\sffamily MICRO-FRONTEND ENGINEERING SUITE}
\fancyhead[R]{\footnotesize\sffamily\docshort}
\fancyfoot[L]{\footnotesize\color{grey}Engineering edition | October 2026}
\fancyfoot[R]{\footnotesize\thepage}
\renewcommand{\headrulewidth}{0.35pt}
\renewcommand{\arraystretch}{1.2}
\setlength{\emergencystretch}{2em}
\lstdefinestyle{eng}{basicstyle=\ttfamily\footnotesize,keywordstyle=\color{navy}\bfseries,commentstyle=\color{grey},stringstyle=\color{green},breaklines=true,breakatwhitespace=false,columns=fullflexible,keepspaces=true,showstringspaces=false,frame=single,rulecolor=\color{border},backgroundcolor=\color{pale},xleftmargin=5pt,xrightmargin=5pt,aboveskip=10pt,belowskip=10pt,captionpos=b}
\lstdefinelanguage{JavaScript}{morekeywords={const,let,var,require,return,async,await,import,export,default,function,try,catch,new,true,false,undefined,module},sensitive=true,morecomment=[l]//,morestring=[b]',morestring=[b]"}
\lstset{style=eng}
\newtcolorbox{sourcebox}{colback=pale,colframe=border,arc=2pt,boxrule=0.6pt,left=8pt,right=8pt,top=7pt,bottom=7pt}
\newtcolorbox{engineerbox}{colback=softgreen,colframe=green!45,arc=2pt,boxrule=0.6pt,left=8pt,right=8pt,top=7pt,bottom=7pt}
\newcommand{\based}[1]{\begin{sourcebox}\textbf{BOOK-GROUNDED CONCEPT}\\#1\end{sourcebox}}
\newcommand{\added}[1]{\begin{engineerbox}\textbf{ENGINEERING EXTENSION}\\#1\end{engineerbox}}
\newcommand{\term}[1]{\textbf{\color{navy}#1}}
\newcommand{\criterion}[1]{\textbf{Acceptance:} #1}
\newcommand{\book}{\emph{Building Micro-Frontends}, 2nd ed.}
\newcommand{\refs}{\chapter*{References and Provenance}\addcontentsline{toc}{chapter}{References and Provenance}
\begin{enumerate}[leftmargin=*]
\item \textbf{Primary source:} Luca Mezzalira, \emph{Building Micro-Frontends: Distributed Systems for the Frontend}, 2nd ed., O'Reilly Media, published November 2025 (copyright 2026), ISBN 978-1-098-17078-3. Chapter and printed-page references are included throughout. This suite paraphrases and synthesizes concepts; it does not reproduce the book.
\item \textbf{Engineering extension:} Webpack Module Federation reference: \url{https://webpack.js.org/plugins/module-federation-plugin}.
\item \textbf{Engineering extension:} Next.js multi-zones: \url{https://nextjs.org/docs/app/guides/multi-zones}. Details are version-sensitive and must be checked against the deployed framework version.
\item \textbf{Engineering extension:} OWASP Application Security Verification Standard 5.0: \url{https://owasp.org/www-project-application-security-verification-standard/}.
\item \textbf{Engineering extension:} GitHub Actions secure use reference: \url{https://docs.github.com/en/actions/reference/security/secure-use}.
\end{enumerate}}

\titleclass{\chapter}{straight}\titleformat{\chapter}[hang]{\huge\sffamily\bfseries\color{navy}}{\thechapter}{.7em}{}\titlespacing*{\chapter}{0pt}{24pt}{13pt}
\newcommand{\docshort}{DECISION FRAMEWORK}
\begin{document}
\begin{titlepage}\thispagestyle{empty}
\vspace*{1.1in}{\Huge\sffamily\bfseries\color{navy} Architecture Decision Framework\par}
\vspace{18pt}{\Large\sffamily\color{blue} Candidate selection, quality-attribute trade-offs, governance, and ADR templates\par}
\vspace{27pt}\noindent\color{blue}\rule{0.75\textwidth}{1.4pt}
\vspace{23pt}\par{\large\sffamily A companion to Luca Mezzalira's \emph{Building Micro-Frontends}, Second Edition}\par
\vfill\noindent\textbf{Document series:} Micro-Frontend Software Architecture and Engineering\\
\textbf{Version:} 1.0 \quad \textbf{Prepared:} October 2026\\
\textbf{Production format:} pdfLaTeX-compatible standalone source\\
\textbf{Scope:} Original synthesis plus explicitly identified engineering extensions
\end{titlepage}
\pagenumbering{roman}\tableofcontents\clearpage\pagenumbering{arabic}

\chapter{Decision Inputs and Non-Goals}
\based{The chapter 2 four-decision framework is the foundation. Chapters 3 and 13 add architecture-characteristic trade-offs and alignment between organization and architecture. This document adds a formal evaluation method; its scores are not from the book.}
\section{Context collection worksheet}
Capture five observations before choosing runtime technology: (1) primary user journeys and where they cross business domains; (2) team size, ownership, and handoff frequency; (3) build/test/deploy lead time and release coupling; (4) traffic mix, SEO requirements, devices, and performance baselines; and (5) critical identity, compliance, and availability constraints.
\section{Disqualifying conditions}
Defer decomposition when product-domain boundaries are unclear, there is no budget for shared delivery tooling and observability, only one team delivers all features, or current delivery bottlenecks stem from poor test design rather than codebase boundaries. Improve the modular monolith first, then reassess.
\chapter{Decision Workflow}
\section{Stage 0: Establish the desired change}
Define expected improvements such as fewer blocked releases, lower coordination time, smaller ownership scope, or faster domain-local deploys. Specify measurement windows and countermetrics including JS shipped, total operating cost, accessibility defects, and incident recovery time.
\section{Stage 1: Identify stable vertical slices}
A good candidate has a coherent business outcome, owner, route or view scope, independently testable behavior, few synchronous cross-domain requirements, and a stable backend contract. If a candidate is only a shared button or form field, keep it a component.
\section{Stage 2: Compare composition approaches}
\begin{longtable}{p{.27\linewidth}p{.33\linewidth}p{.31\linewidth}}\toprule Constraint & Preferred starting candidate & Watch out for \\\midrule
Independent route domains & URL/vertical split or multi-zones & Hard navigation, session continuity\\
Dense cross-domain dashboard & Client-side horizontal composition & Runtime drift and startup cost\\
SEO and initial HTML critical & Route-based SSR or controlled server fragments & Origin latency and personalization cache\\
Different trust levels & Separate documents/origins as needed & Iframe communication and identity\\
High client framework diversity & Route/zone split & Shared design semantics\\
Low platform maturity & Modular frontend monolith & Revisit independent deployment later\\\bottomrule\end{longtable}
\section{Stage 3: Specify routing and communication}
Define owner for each URL prefix; links or browser history semantics; auth and session propagation; event names and schemas; remote failure paths; access-control checks; data ownership; and contract versioning. Avoid bidirectional remote imports. Draw data-flow boundaries separately from UI composition boundaries.
\section{Stage 4: Validate with a bounded experiment}
Prototype exactly one high-value capability; run happy path, direct-link, refresh, logout, remote outage, previous-release compatibility, degraded network, and rollback tests. A proof of concept without fault injection cannot demonstrate operational viability.
\chapter{Weighted Trade-Off Analysis}
\added{The scorecard below is an illustrative engineering instrument. The weights and scores must be agreed by product, platform, frontend, security, and operations stakeholders; they are not universal empirical values.}
\section{Example scoring model}
For candidate architecture $a$, let $S(a)=\sum_{i=1}^{n}w_i r_{ia}$, where weights $w_i\geq0$ sum to $1$, and ordinal ratings $r_{ia}\in\{1,2,3,4,5\}$ are higher when fit is better. Use absolute failure gates before computing rank.
\begin{center}\small\begin{tabular}{p{.32\linewidth}rccc}\toprule Attribute & Weight & Modular SPA & Route zones & Runtime composition\\\midrule
Team independence & .25 & 2 & 5 & 5\\
User-perceived performance & .20 & 4 & 4 & 2\\
Operational simplicity & .20 & 5 & 3 & 2\\
Cross-domain experience & .15 & 4 & 3 & 5\\
Deployment isolation & .10 & 2 & 5 & 4\\
Existing skills/tooling & .10 & 5 & 4 & 3\\\midrule
Weighted score & 1.00 & 3.60 & 4.00 & 3.50\\\bottomrule\end{tabular}\end{center}
Scores show the route-zone candidate ahead under \emph{these assumptions}. Re-run with a sensitivity test, e.g., increase performance weight at the expense of autonomy, and identify the reversal threshold. An integer score does not substitute for a load test or threat model.
\section{Hard gates}
\begin{enumerate}\item Server authorization for every protected API and workflow.
\item Browser security model: trusted remote origins, CSP strategy, cookie/session design.
\item Resource and performance budgets that keep the core user path usable.
\item Detectable failures and tested deployment rollback.
\item Design and accessibility continuity for core journeys.
\end{enumerate}
\chapter{Boundary and Dependency Heuristics}
\section{Bounded-context rubric}
Assign each proposed domain a 0--2 score for: business cohesion, explicit owner, independent release, data-contract stability, UX separability, and failure containment (maximum 12). A low score flags the need for a redesign workshop; it does not mathematically prove the boundary is wrong. Review high-coupling edges with a context map.
\section{Integration contract}
Every contract specifies \textbf{producer, consumer, schema, version, compatibility window, error model, security assumption, test fixture, observability attributes, and retirement date}. Data exchange should remain minimal: a product ID or checkout token is usually safer to share than an entire mutable product or user object. Never treat UI possession of an identifier as authorization.
\chapter{Architecture Decision Records}
\section{ADR template}
\begin{sourcebox}
\textbf{ADR-XXX -- [Decision Name]}\par
Status: proposed / accepted / superseded\par
Context: [measurable problem; affected domains]\par
Decision: [boundary, composition, routing, communication]\par
Alternatives: [at least two, including no-decomposition]\par
Consequences: [performance, deployment, security, team impact]\par
Validation: [test, metric, benchmark, migration pilot]\par
Ownership: [approvers, maintainer, review trigger]\par
Rollback: [how to reverse or isolate the decision]
\end{sourcebox}
\section{Worked example: customer account boundary}
\textbf{Context:} Account changes require separate audit and release control; frequent catalog enhancements should not wait for identity workflows. \textbf{Decision:} Extract \texttt{/account/*} into a route-owned application using an existing edge router, backend session validation, and a shared visual token package. \textbf{Rejected:} Build a widget-based horizontal federation inside the catalog page, because the account flow is navigational and does not demand same-view composition. \textbf{Validation:} Independent deploy, login/logout, deep link, session expiry, accessibility, and rollback. \textbf{Risk:} Cross-zone navigation reloads and identity cookie scope. \textbf{Trigger to revisit:} A validated requirement for dense same-view account widgets.
\chapter{Governance and Decision Review}
\based{Chapter 13, pp. 343--367, addresses team responsibilities, communication, RFCs, ADRs, and decentralized organizations.}
\begin{tabularx}{.96\linewidth}{p{.23\linewidth}X}\toprule Accountable party & Responsibilities \\\midrule
Domain teams & Business contracts, delivery, local error handling, on-call, security fixes\\
Frontend platform & Shell/discovery contracts, build blueprints, dependency policy, shared telemetry\\
Security/AppSec & Threat-model patterns, control verification, exception process, supply-chain requirements\\
Architecture forum & Cross-domain ADRs, performance guardrails, compatibility escalation\\
Product/UX & User journey integrity, accessibility consistency, prioritization\\\bottomrule\end{tabularx}
\section{Review cadence and triggers}
Review ADRs on major framework changes, ownership changes, incidents, material performance degradation, or newly coupled releases. Monitor \term{independent deploy rate}, \term{change failure rate}, \term{lead time}, \term{restoration time}, and domain-level user outcomes. These are supporting measures, not reasons to game deployment counts.
\chapter{Architecture Review Checklist}
\begin{itemize}[label=$\square$]\item Business objective and baseline metric present.
\item Each domain has an accountable owner and lifecycle.
\item All four architectural decisions explicitly documented.
\item Modular monolith alternative examined.
\item Sensitive data and trust boundaries identified.
\item Dependency/shared-state graph contains no unexplained cycles.
\item Failure modes tested under degraded network and remote outages.
\item Direct navigation, accessibility, and SEO requirements measured.
\item Staged promotion, recovery target, and rollback owner named.
\item Decision review date, triggers, and success thresholds recorded.\end{itemize}

\refs
\end{document}
